\vspace{-6mm}
\section{Introduction}
\vspace{-4mm}
Machine learned models can accelerate the prediction of quantum properties of atomistic systems like molecules by learning approximations of \textit{ab initio} calculations ~\citep{neural_message_passing_quantum_chemistry, deep_potential_molecular_dynamics, push_limit_of_md_100m, dimenet_pp, nequip,  deep_potential_molecular_dynamics_simulation, se3_wavefunction, gemnet_xl, quantum_scaling}.
%Particularly, graph neural networks (GNNs), which naturally treat the set-like nature of collections of atoms and mimic the topological and geometric nature of atomic interactions, have gained increasing popularity due to their performance. 
%, which is similar to convolutional neural networks (CNNs).
%Neural networks are top performers in existing atomistic benchmarks due to the ease with which inductive biases can be incorporated to exploit the symmetry of training data.
% One main factor leading to the success of neural networks is incorporating inductive biases to exploit the symmetry of training data.
In particular, graph neural networks (GNNs) have gained increasing popularity due to their performance.
By modeling atomistic systems as graphs, GNNs naturally treat the set-like nature of collections of atoms, encode the interaction between atoms in node features and update the features by passing messages between nodes.
One factor contributing to the success of neural networks is the ability to incorporate inductive biases that exploit the symmetry of data.
% One factor contributing to the success of neural networks is the ability to encode inductive biases (e.g. the symmetry of data) into the functional form of the neural network.
%\revision{One factor contributing to the success of neural networks is the ability to encode inductive biases into functional spaces.}
Take convolutional neural networks (CNNs) for 2D images as an example: Patterns in images should be recognized regardless of their positions, which motivates the inductive bias of translational equivariance. % in convolutional layers.
As for atomistic graphs, where each atom has its coordinate in 3D Euclidean space, we consider inductive biases related to 3D Euclidean group $E(3)$, which include equivariance to 3D translation, 3D rotation, and inversion.
Concretely, some properties like energy of an atomistic system should be constant regardless of how we shift the system; others like force should be rotated accordingly if we rotate the system.
To incorporate these inductive biases, equivariant and invariant neural networks have been proposed.
% ============================================
% ============================================
%Graph neural networks (GNNs) have gained increasing popularity for building models of quantum properties of atomic systems as they naturally treat the set-like nature of collections of atoms and mimic the topological and geometric nature of atomic interactions. 
% e.g. atomic ``node'' features are iterative updated through message-passing along ``edges'' defined by bonds or geometrically adjacent atoms.
%Additionally, each atom can be assigned coordinates in 3D Euclidean space, but the choice of coordinate system is fundamentally arbitrary; we can always change between coordinate systems via 3D translations, rotations and inversion, i.e. the symmetries 3D Euclidean space. These symmetries impact predicted properties in different ways:
% In this work, we additionally treat the fact that each atom has coordinates in 3D Euclidean space and consider inductive biases related to 3D Euclidean group $E(3)$, which include equivariance to 3D translation, 3D rotation, and inversion.
% Concretely, 
%Some properties like the energy of an atomistic system are constant regardless of how the coordinate system changes; others like force
% change (e.g. rotate) with the change in coordinate system.
%e.g. rotate with rotation of the coordinate system
% Some properties like energy of an atomistic system are constant regardless of how we change the coordinate system; others like force rotate accordingly if we rotate the system.
%To incorporate the inductive bias of 3D Euclidean symmetry, equivariant and invariant neural networks have been proposed.
% ============================================
% ============================================
The former leverages geometric tensors like vectors for equivariant node features~\citep{tfn, 3dsteerable, kondor2018clebsch, se3_transformer, nequip, segnn, allergo}, and the latter augments graphs with invariant information such as distances and angles extracted from 3D graphs~\citep{schnet, dimenet, dimenet_pp, spherenet, gemnet}.

\vspace{-1.5mm}
A parallel line of research focuses on applying Transformer networks~\citep{transformer} to other domains like computer vision~\citep{detr, vit, deit} and graph~\citep{generalization_transformer_graphs, spectral_attention, graphormer, graphormer_3d} and has demonstrated widespread success.
However, as Transformers were developed for sequence data~\citep{bert, wav2vec2, gpt3}, it is crucial to incorporate domain-related inductive biases. % when generalizing Transformer.
For example, Vision Transformer~\citep{vit} shows that adopting a pure Transformer to image classification cannot generalize well and achieves worse results than CNNs when trained on only ImageNet~\citep{imagenet} since it lacks inductive biases like translational invariance.
Note that ImageNet contains over 1.28M images and the size is already larger than that of many quantum properties prediction datasets~\citep{qm9_2, md17_1, oc20}.
Therefore, this highlights the necessity of including correct inductive biases when applying Transformers to the domain of 3D atomistic graphs.


%In this work, we 
% study how to incorporate $E(3)$-related inductive biases into Transformer and 
%present \textbf{Equiformer}, an equivariant graph neural network utilizing $SE(3)$/$E(3)$-\textbf{equi}variant features built from irreducible representations (irreps) and $E(3)$-\textbf{equi}variant attention mechanisms to combine the inductive bias of E(3) symmetry with  
% and leveraging 
%In this work, we present \textbf{Equiformer}, an equivariant graph neural network utilizing $SE(3)$/$E(3)$-\textbf{equi}variant features built from irreducible representations (irreps) and a novel \textbf{equi}variant attention mechanism to combine the 3D-related inductive bias with the strength of Trans\textbf{former}.
\vspace{-1.5mm}
\revision{Despite their widespread success in various domains, Transformers have yet to perform well across datasets~\citep{se3_transformer, torchmd_net, eqgat} in the domain of 3D atomistic graphs even when relevant inductive biases are incorporated.}
In this work, we \revision{demonstrate that Transformers can generalize well to 3D atomistic graphs} and present \textbf{Equiformer}, an equivariant graph neural network utilizing $SE(3)$/$E(3)$-\textbf{equi}variant features built from irreducible representations (irreps) and a novel attention mechanism to combine the 3D-related inductive bias with the strength of Trans\textbf{former}.
%%%Irreps features encode equivariant information in channels without complicating graph structures.
%%%The simplicity enables us to directly incorporate them into Transformers through replacing original operations with equivariant counterparts and introducing an additional equivariant operation called tensor product.
%Moreover, we propose a novel equivariant graph attention, which considers both content and geometric information such as relative position. % contained in irreps features.
%%%Moreover, we propose a novel attention mechanism called equivariant graph attention, which considers both content and geometric information such as relative position. % contained in irreps features.
%Equivariant graph attention improves upon typical attention in Transformers by replacing dot product attention with theoretically stronger
%a mathematically more expressive
%multi-layer perceptron attention and including non-linear message passing.
%%%The proposed attention improves upon typical attention in Transformers by replacing dot product attention with theoretically stronger multi-layer perceptron attention and including non-linear message passing.
%With these innovations, Equiformer demonstrates the possibility of generalizing Transformer-like networks to 3D atomistic graphs and achieves competitive results on two quantum properties prediction datasets, QM9~\cite{qm9_2, qm9_1} and OC20~\cite{oc20}.
%%%With these innovations, Equiformer demonstrates the possibility of generalizing Transformers to 3D atomistic graphs and achieves competitive results on QM9~\citep{qm9_2, qm9_1}, MD17~\citep{md17_1, md17_2, md17_3}, revised MD17~\citep{revised_md17} and OC20~\citep{oc20} datasets.
%%%With these innovations, Equiformer achieves competitive results on QM9~\citep{qm9_2, qm9_1}, MD17~\citep{md17_1, md17_2, md17_3}, revised MD17~\citep{revised_md17} and OC20~\citep{oc20} datasets.
\revision{First, we propose a simple and effective architecture, Equiformer with dot product attention and linear message passing, by only replacing original operations in Transformers with their equivariant counterparts and including tensor products.}
\revision{Using equivariant operations enables encoding equivariant information in channels of irreps features without complicating graph structures.}
\revision{With minimal modifications to Transformers, this architecture has already achieved strong empirical results (Index 3 in Table~\ref{tab:qm9_ablation} and~\ref{tab:oc20_ablation}).}
\revision{Second, we propose a novel attention mechanism called equivariant graph attention, which improves upon typical attention in Transformers through replacing dot product attention with multi-layer perceptron attention and including non-linear message passing.}
%With these innovations, Equiformer achieves competitive results on QM9~\citep{qm9_2, qm9_1}, MD17~\citep{md17_1, md17_2, md17_3} and OC20~\citep{oc20} datasets.
\revision{Combining these two innovations,} Equiformer \revision{(Index 1 in Table~\ref{tab:qm9_ablation} and~\ref{tab:oc20_ablation})} achieves competitive results on QM9~\citep{qm9_2, qm9_1}, MD17~\citep{md17_1, md17_2, md17_3} and OC20~\citep{oc20} datasets.
%For QM9, compared to models trained with the same data partition, Equiformer achieves the best results on 11 out of 12 regression tasks.
%\revision{For QM9, Equiformer achieves overall better results than previous models across $12$ regression tasks.}
%\revision{For MD17 and , Equiformer achieves overall better results than previous models across $12$ regression tasks.}
%For OC20, under the setting of training with IS2RE data and optionally IS2RS data, Equiformer improves upon state-of-the-art models.
%\todo{[Detailed examples of MD17, revised MD17.]}
For QM9 and MD17, Equiformer achieves overall better results across all tasks or all molecules compared to previous models like NequIP~\citep{nequip} and TorchMD-NET~\citep{torchmd_net}.
%%%For OC20, under the setting of training with IS2RE data and optionally IS2RS data, Equiformer improves upon state-of-the-art models such as SEGNN~\citep{segnn} and Graphormer~\citep{graphormer_3d}.
For OC20, when trained with IS2RE data and optionally IS2RS data, Equiformer improves upon state-of-the-art models such as SEGNN~\citep{segnn} and Graphormer~\citep{graphormer_3d}.
Particularly, as of the submission of this work, Equiformer achieves the best IS2RE result when only IS2RE and IS2RS data are used and improves training time by $2.3 \times$ to $15.5 \times$ compared to previous models.

%The contributions of this work are as follows:
%\vspace{-1.5mm}
%\begin{enumerate}
%    \item We show that through utilizing equivariant irreps features and equivariant operations like tensor product, we can incorporate $E(3)$-related inductive biases into Transformer. \vspace{-0.3mm}
%    \item We propose equivariant graph attention, which considers both content and geometric information in irreps features and improves upon typical attention in Transformer through utilizing multi-layer perceptron attention and non-linear message passing. \vspace{-0.3mm}
%    \item With the proposed methods, Equiformer improves upon state-of-the-art results on QM9 and OC20 datasets. 
%\end{enumerate}
    
    